\section{Reflection、オープンウェイトモデル公開準備をAxiosが報道}
Axiosは、NVIDIA出資のReflectionが中国の上位オープンウェイトと競争し得るモデルを近く公開すると報じた。窓内のX投稿はこの報道の要約が中心で、モデル名、パラメータ数、ライセンス、ベンチマーク、公開日は未確認。NebiusやSpaceXとの計算資源契約にも言及されるが、未名ソースを含む報道段階である。
\dailyxsource{Andrew Curran (@AndrewCurran\_)}{https://x.com/AndrewCurran_/status/2106846163534241912}

\section{Geminiアプリ、無料版は10月9日からFlash-Liteのみとの報道}
窓内の日本語投稿は、個人のGemini無料アカウントが10月9日からFlash-Liteのみとなり、FlashとProを選べなくなると伝えた。AI PlusはProを失い、AI Pro / Ultraは3系統を維持するとの整理もある。Grok収集ではGoogleサポート文書の原文URLを確認していない。
\dailyxsource{古澤久志｜AI×経営実践 (@h\_www)}{https://x.com/h_www/status/2106856501306671113}

\section{米政府「Super Intelligence Force」発足をFox Newsが投稿}
Fox Newsは、Trumpが米国のAI優位を確保するため「Super Intelligence Force」を発表したと投稿した。関連まとめでは参加機関や責任者にも触れるが、Grok収集では大統領令や省庁発表の原文URLを取得していない。任務範囲、法的権限、予算は未確認である。
\dailyxsource{Fox News (@FoxNews)}{https://x.com/FoxNews/status/2106724233510629751}

\section{SpaceXAIをSpaceXSIへ改称するとの報道}
Watcher.Guruは、Elon MuskがSpaceXAIをSpaceXSIへ改称すると投稿した。ForbesもMuskがX上で改称に応じたと報じているが、Grok収集ではMusk本人の該当投稿IDを取得できなかった。記事時点でサイトやアカウント名は未変更で、実施時期も未提示とされる。
\dailyxsource{Watcher.Guru (@WatcherGuru)}{https://x.com/WatcherGuru/status/2106752902203163081}

\section{Grok 4.7、Artificial Analysis Cyber Indexで1位との投稿}
X Freezeは、Grok 4.7 xHighがArtificial AnalysisのCyber Indexで1位になったと画像付きで投稿し、Muskも同趣旨で引用した。Grok収集ではリーダーボードURL、モデル設定、xHighの定義、比較対象のバージョンや再現条件を確認していない。
\dailyxsource{X Freeze (@XFreeze)}{https://x.com/XFreeze/status/2106773576258851087}

\section{Kolibri-1を窓内投稿が紹介、公式発表は観測窓外}
窓内の投稿はAleph AlphaのKolibri-1について、78B総量・約3.5B activeのMoE、最大100万トークン、独英対応、Apache 2.0などを紹介した。公式投稿は今回の観測窓外で、ベンチマーク値も同社自己評価として扱われている。
\dailyxsource{Sinda | SKatalyst AI (@SKatalystAI)}{https://x.com/SKatalystAI/status/2106844692503466098}

\section{GoogleがOSSバグバウンティを停止したとの二次投稿}
Applitankは、AI生成投稿の急増を受けGoogleがオープンソース向けバグバウンティを凍結したとTechCrunch経由で投稿した。別の日次まとめでも関連情報が述べられるが、いずれも二次情報である。Googleの公式声明やVRPページ更新はGrok収集では未確認。
\dailyxsource{Applitank (@applitank)}{https://x.com/applitank/status/2106852257950228786}

\section{ELYZAが「RSI Research」を始めたとの日本語まとめ}
日本語AIニュースまとめは、ELYZAが「ELYZA RSI Research」を設立し、自律改善実験で精度をほぼ維持しながら推論コストを67\%削減したと紹介した。投稿者自身も現段階はRSIそのものではなく前段階と留保している。公式発表URLや実験条件は未確認。
\dailyxsource{BB (@Bb17350260Bb)}{https://x.com/Bb17350260Bb/status/2106857625480757421}

\section{Meta MuseをESP32 / Raspberry PiへつなぐSDK公開とのまとめ}
同じ日本語まとめは、MetaがMuseをESP32やRaspberry Piへ接続するSDKとファームウェアをオープンソース化したと伝えた。Grok収集ではMeta公式投稿やリポジトリURLを確認しておらず、対応ボード、ライセンス、認証方式も不明である。
\dailyxsource{BB (@Bb17350260Bb)}{https://x.com/Bb17350260Bb/status/2106857625480757421}

\section{ローカルQwenとクラウド課金のコスト比較投稿}
Digital Bathは、Ryzen AI 9搭載64GBミニPCで35B Qwenを推論し、1日600万トークン相当の処理をクラウド利用料と比較したという記事を紹介した。投稿では13--17 tok/sなどの数値を挙げるが、原記事、正確なモデルID、量子化、電力、品質の等価性は未確認。
\dailyxsource{Digital Bath (@FullDigitalBath)}{https://x.com/FullDigitalBath/status/2106863800632873147}

\section{次のオープンウェイトモデルをめぐる噂リスト}
Chubbyは、次のオープンウェイトの「rumor watchlist」としてQwen 4、Kimi K3.1、次のGLM、DeepSeek V4.2を挙げた。10月から11月という時期も推測と明記しており、公開事実ではない。
\dailyxsource{Chubby (@kimmonismus)}{https://x.com/kimmonismus/status/2106863682663882897}

\section{Karpathy、AI関心層の新規参入について所感}
Andrej Karpathyは、現在AIに注目している人の99\%以上は参加から1年未満で、それ以前から関わる人には状況が混乱して見えるという所感を投稿した。数字を裏付ける定量調査は示されておらず、本人のメンタルモデルとして読むべき内容である。
\dailyxsource{Andrej Karpathy (@karpathy)}{https://x.com/karpathy/status/2106806571321966793}

\section{LeCun、計画できる系はLLMではないと返信}
Yann LeCunは、計画し、行動結果を予測し、複数の帰結を比較できるなら世界モデルを持っており、それはLLMではないと返信した。新しい論文やモデルの発表ではなく、研究観に関する持論として観測された投稿である。
\dailyxsource{Yann LeCun (@ylecun)}{https://x.com/ylecun/status/2106811742319698152}

\section{BitBender、難読化バイナリ前処理サービスをソフトローンチ}
Aftermath Labsは、難読化されたバイナリをLLM解析へ渡す前段の処理を行うBitBenderをソフトローンチし、ベータ利用者を募集した。Grok収集では製品ページや第三者による効果検証は確認していない。
\dailyxsource{Aftermath Labs (@aftermath\_labs)}{https://x.com/aftermath_labs/status/2106854872020824361}
